% Abstrat of your communication to the
% ACA2018 Conference (Santiago de Compostela, 2018).
% You must use LaTeX format.


%====================================================
% IMPORTANT: DO NOT MODIFY THE FOLLOWING LINES
%====================================================

\documentclass{article}

\usepackage[T1]{fontenc}
\usepackage[utf8]{inputenc}
\usepackage{amssymb,amsthm,amsmath,amsfonts,latexsym}
\usepackage{graphicx}
\usepackage{hyperref}
\usepackage{times}

\setlength{\oddsidemargin}{1.46cm}
\setlength{\textwidth}{13cm}
\setlength{\topmargin}{0.0cm}
\setlength{\textheight}{20.5cm}

\makeatletter
\newcommand{\TITLE}[2][]{
  \begin{center} \Large \bfseries #2%
  \def\@temp{#1}\ifx\@temp\@empty\relax\else\footnote{#1}\fi\end{center}}

\newcommand{\AUTHOR}[2][]{\textbf{#2}%
  \def\@temp{#1}\ifx\@temp\@empty\relax\else\textsuperscript{#1}\fi}

\newcommand{\ADDRESS}[2]{\noindent
  \textsuperscript{#1}\parbox[t]{0.9\textwidth}{#2} \vspace{6pt}}

\makeatother

\newcommand{\KEYWORDS}[1]{\paragraph{Keywords:} #1}
\newcommand{\MSCLASSIFICATION}[1]{\paragraph{Mathematics Subject Classification 2010:} #1}

\renewcommand{\thefootnote}{\fnsymbol{footnote}}
\setcounter{footnote}{0}
\pagestyle{empty}

\newcommand{\HEADING}{
{\noindent \small Applications of Computer Algebra -- ACA2018 \\
Santiago de Compostela, June 18--22, 2018} \vspace{10pt}}

%===============================================
% DO NOT MODIFY THE ABOVE LINES
%===============================================

\begin{document} 

\HEADING


\TITLE{Automatic generation of diagrammatic subway maps for any date
       with Maple}

% If you need to thank someone next to the title, please use the next line (and delete the previous one)

% \TITLE[thanks]{Title of talk}

% Write the author names. Identify each one of them with a number
% in order to give their own affiliations later on.
\begin{center}
 \AUTHOR[1]{Alberto Almech},
 \AUTHOR[2]{Eugenio Roanes-Lozano}
% IF two or more authors have the same affiliation, associate to them the same number.
% For instance, if there are three authors and the first a the second ones have the same affiliation,
% write as follows:
 % \AUTHOR[1]{Third Author}
\end{center}


% Write now the abstract of your communication.

The second author was one of the authors of 
a computer package written in Maple that could automatically 
generate railway maps of a network for any date. This package was
presented at ACA'2008 and its design and implementation is described in [1]. 
Each section of the network was coloured accordingly to its characteristics
(single / double track, electrified / non electrified,
opened / closed / greenway,...). 
The position of the nodes (stations, junctions,...) was obtained from 
a list of geographical coordinates.

The work presented here deals with a similar although not identical 
case: subway networks are treated as graphs with the help of a computer
algebra system in order to obtain the diagrammatic map for any date. 

Most metro network plans follow
more or less closely the ideas introduced by Harry Beck in his  
diagrammatic design of London subway map (the distances between stations
and geographic orientation of the lines don't have to be respected, as
the clarity and the number of stations between two stations is
the key information to be visualized). 

Therefore allocating nodes is far simpler, and we have decided to 
manually allocate the stations on a predefined grid.

The situation is also simpler because all lines
are double track and electrified. For instance in Madrid subway 
there are minor differences between lines, such as the kind of catenary 
(classic or rigid), the gauge (narrow / broad),... that will not be 
considered here. Each node and edge of the graph has dates 
associated: inauguration date / closure date --the latter if applies.

The package takes advantage of the simplifications w.r.t. [1]
mentioned above
and the features of \textit{Maple's} \textit{Networks} package.
This way the approach, although general, can be implemented in 
relatively few lines of code. 

We know of no other similar works.

The work is illustrated with the case of Madrid subway network,
one of the biggest ones in the world.

% Write the keywords separated by commas

\KEYWORDS {Graph theory, Network models, Diagrammatic maps, Subways}

% Write the MSC2010 codes corresponding to your communication

\MSCLASSIFICATION{68R10, 90B10, 90C35}
                  % Graph theory (including graph drawing)
                  %	Network models, deterministic
                  % Programming involving graphs or networks
%
% Next include the references you need.
% If your abstract includes no references at all, please delete all the lines
% from \begin{thebibliography} to \end{thebibliography}.

% Please, do not use more than 9 references.

\begin{thebibliography}{9}

 %  Model for references to articles
  \bibitem{label1} \textsc{E. Roanes-Lozano, A. Mart\'{\i}nez-Zarzuelo, 
                           A. Garc\'{\i}a-\'Alvarez, 
                           M. J. Wester, E. Roanes-Mac\'{\i}as}
                    Automatically Obtaining Railway Maps from a Set of Historical 
                    Events
                    \emph{ RACSAM (Revista de la Real Academia de Ciencias Exactas,      
                     F\'{\i}sicas y Naturales, Serie A, Matem\'aticas)}
                    \textbf{105}(1), 149--165 (2011).
                    DOI 10.1007/s13398-011-0010-1
\end{thebibliography}

\vspace{1cm}

% Full Affiliation of all authors
\ADDRESS{1}{
%Department1  \\
            Facultad de CC. Matem\'aticas, Universidad Complutense de Madrid\\ 
            Plaza de Ciencias s/n, 28040-Madrid, Spain \\
\texttt{albermech@gmail.com}


\vspace{.5cm}

\ADDRESS{2}{
            Instituto de Matem\'atica Interdisciplinar \& \\
            Depto. de \'Algebra, Geometr\'{\i}a y Topolog\'{\i}a\\
            Facultad de Educaci\'on, Universidad Complutense de Madrid\\ 
            c/ Rector Royo Villanova s/n, 28040-Madrid, Spain \\
\texttt{eroanes@mat.ucm.es}
}
}


\end{document}
